% META: {"id":"TCOMPL-ECH-CO-002","chapitre":"TCOMPL-ECHANTILLONNAGE","type_objet":"corrige","exercice_ref":"TCOMPL-ECH-EX-002","capacites_codes":["C2"],"sources_inspiration":[],"status":"approved","origin":"generated","status_history":["generated","audited","accepted","approved"],"release_acceptance":"RELEASE_OWNER_BATCH_ACCEPTANCE","acceptance_closure_digest":"sha256:9b3ccf9a81c5520fbb7e03b7b2d3e7bf2057a02834b6d908bf3be5f49f3b553f","acceptance_receipt_digest":"sha256:4fad86f0282e0fa4e0419cca1ad6379ae7af5a280c04a4a90880f90a1c044242"}
% BEGIN-VERIFY
% from sympy import binomial
% assert binomial(4, 0) == 1
% assert binomial(4, 1) == 4
% assert binomial(4, 2) == 6
% assert binomial(4, 3) == 4
% assert binomial(4, 4) == 1
% END-VERIFY
\begin{corrige}{TCOMPL-ECH-EX-002}
En additionnant deux à deux les coefficients de la ligne $n=3$ ($1,3,3,1$), en complétant par des $1$ aux extrémités :
\[
  \binom{4}{0}=1,\quad \binom{4}{1}=1+3=4,\quad \binom{4}{2}=3+3=6,\quad \binom{4}{3}=3+1=4,\quad \binom{4}{4}=1.
\]
\end{corrige}
